\documentclass[11pt]{article}

\usepackage[T1]{fontenc}
\usepackage{lmodern}
\usepackage{amsmath,amssymb,amsthm}
\usepackage[margin=1in]{geometry}
\usepackage{microtype}

\newtheorem{theorem}{Theorem}
\newtheorem{lemma}[theorem]{Lemma}
\newtheorem{proposition}[theorem]{Proposition}

\newcommand{\ND}{N_D}
\newcommand{\Rc}{\mathcal R_{\rm c}}
\newcommand{\Dc}{\mathcal D_{20}}
\newcommand{\Rrr}{\mathcal R_{\rm rr}}

\title{Purely Analytic Closure of the Final Spherical-Shell Residual}
\author{}
\date{}

\begin{document}
\maketitle

\begin{abstract}
This note assembles the all-frequency retained-remainder estimate after the
exact owner subtraction. It proves that the final spherical-shell residual
is empty without invoking an interval certificate, an executable truth
table, an aggregate tail argument, or a frequency split.  The structural,
angular-block, and scalar-positivity modules cover one closed superset of the
entire residual.
\end{abstract}

\section{Notation and analytic inputs}

Let
\[
 A_{\rho,1}:=\{x\in\mathbb R^3:\rho<|x|<1\},
 \qquad
 W(\rho,K):=\frac{2(1-\rho^3)K^3}{9\pi},
\]
and let \(\ND(A_{\rho,1},K^2)\) count Dirichlet eigenvalues strictly below
\(K^2\). Put
\[
 \rho_c:=\frac1{1+2\pi},
 \qquad
 k_{11}(\rho):=
 \sqrt{\frac{\pi^2}{(1-\rho)^2}+132}.
\]

The preceding analytic owner lemmas and the exact set subtraction give
\begin{equation}
 \Dc=
 \left\{(\rho,K):
 \rho_c\leq\rho<\frac{39}{50},\quad
 k_{11}(\rho)<K<U(\rho)\right\}.
 \label{eq:D20}
\end{equation}
On this interval \(U(\rho)=K_0(\rho)\), although the formula for the upper
wall will not be needed below.

We use the shell action
\[
 G_a(z):=\frac1\pi
 \left(\sqrt{a^2-z^2}-z\arccos\frac za\right)
 \quad(0\leq z<a),
 \qquad G_a(z):=0\quad(z\geq a),
\]
\[
 A_{\rho,K}(z):=G_K(z)-G_{\rho K}(z),
\]
and the strict lattice proxy
\begin{equation}
 P(\rho,K):=
 \sum_{\substack{\nu=\ell+1/2<K\\\ell\in\mathbb N_0}}
 2\nu\,
 \#\left\{n\in\mathbb N:
 n<A_{\rho,K}(\nu)+\frac14\right\}.
 \label{eq:P}
\end{equation}
The analytic Bessel-phase reduction proved earlier in the main argument is
\begin{equation}
 \ND(A_{\rho,1},K^2)\leq P(\rho,K).
 \label{eq:phase}
\end{equation}

Put
\[
 M_c(\rho,K):=\frac{(1-\rho^2)K^2}{6}-\frac12.
\]
Let
\begin{equation}
 N:=\#\left\{n\geq1:n-\frac14<\frac{(1-\rho)K}{\pi}\right\}.
 \label{eq:N}
\end{equation}
Parts VII--VIII of the detailed analytic supplement combine the retained
radial remainder with one disjoint left block per active layer and give the
single direct estimate
\begin{equation}
 W(\rho,K)-P(\rho,K)
 >\mathcal H(\rho,K)-\frac{(1-\rho^2)K^2}{6}+\frac N2.
 \label{eq:direct-gap-input}
\end{equation}
On
\(\rho_c\leq\rho\leq39/50\) and \(K\geq k_{11}(\rho)\),
one has \((1-\rho)K/\pi>1\), hence \(N\geq1\).  Part~IX gives on the same
unbounded domain
\begin{equation}
 \mathcal H(\rho,K)-M_c(\rho,K)>0.
 \label{eq:scalar-input}
\end{equation}
The all-frequency scope in \eqref{eq:scalar-input} is part of that module:
its scalar gap has positive \(K\)-derivative whenever \(K>12\), one has
\(k_{11}(\rho)>241/20>12\), and the gap is strictly positive on the base
curve \(K=k_{11}(\rho)\).  Thus the proof integrates upward from the base
curve to every finite \(K\geq k_{11}(\rho)\); it uses no upper-frequency
bound.

\section{All-frequency retained-remainder closure}

\begin{theorem}[All-frequency middle-ratio theorem]
\label{thm:all-frequency}
For
\[
 \rho_c\leq\rho\leq\frac{39}{50},
 \qquad
 K\geq k_{11}(\rho),
\]
one has
\[
 \ND(A_{\rho,1},K^2)<W(\rho,K).
\]
\end{theorem}

\begin{proof}
Since \(N\geq1\), equations \eqref{eq:direct-gap-input} and
\eqref{eq:scalar-input} give
\[
 W-P>\mathcal H-M_c>0.
\]
The phase reduction \eqref{eq:phase} completes the proof.  Strictness includes
both closed ratio faces, the base face \(K=k_{11}(\rho)\), and every common
radial--angular proxy wall.
\end{proof}

\section{Direct ownership of the exact residual}

\begin{proposition}[Analytic closure of \(\mathcal D_{20}\)]
\label{prop:closure}
Define the closed all-frequency theorem domain
\[
 \Rrr:=\left\{(\rho,K):
 \rho_c\leq\rho\leq\frac{39}{50},\quad
 K\geq k_{11}(\rho)\right\}.
\]
Then \(\Dc\subset\Rrr\), the strict P\'olya inequality holds throughout
\(\Dc\), and the successor residual is empty.
\end{proposition}

\begin{proof}
If \((\rho,K)\in\Dc\), then \eqref{eq:D20} gives
\[
 \rho_c\leq\rho<\frac{39}{50},
 \qquad
 k_{11}(\rho)<K<U(\rho).
\]
After weakening the strict upper ratio and lower frequency inequalities,
these conditions place \((\rho,K)\) in \(\Rrr\).  Therefore
Theorem~\ref{thm:all-frequency} proves every point of \(\Dc\), and
\[
 \Dc\setminus\Rrr=\varnothing.
\]
The already owned faces \(\rho=39/50\), \(K=k_{11}(\rho)\), and
\(K=U(\rho)\) remain outside \(\Dc\), even though the stronger theorem
also covers the first two.  The face \(\rho=\rho_c\) remains in \(\Dc\)
whenever both strict frequency inequalities hold.  No artificial
frequency face is present.
\end{proof}

\paragraph{Logical dependency.}
After Parts I--VI identify \(\mathcal D_{20}\), the closure uses only the
structural, angular-block, and scalar-positivity modules in Parts VII--IX.
It does not use an aggregate-tail module, a frequency split, the former
\(10{,}580\)-box compact certificate, the \(1{,}286\)-box aggregate
certificate, or the \(63\)-state executable subtraction table.

\end{document}
